\textbf{Phần 1. Trắc nghiệm lựa chọn}
\setcounterpageref{ex}{0}

\begin{ex} %[5H9V53]
Nếu độ dài cạnh của một hình lập phương gấp lên $3$ lần thì thể tích của hình lập phương đó gấp lên bao nhiêu lần?
\choice
{$9\text{ lần}$}
{\True $27\text{ lần}$}
{$6\text{ lần}$}
{$18\text{ lần}$}
\loigiai{
Thể tích ban đầu: $V = a \times a \times a$.\\
Khi cạnh gấp lên $3$ lần:
$$V' = (3 \times a) \times (3 \times a) \times (3 \times a) = 27 \times (a \times a \times a) = 27 \times V.$$
Vậy thể tích gấp lên $27$ lần.\\
\textbf{Đáp án của câu hỏi là B}
}
\end{ex}

\begin{ex} %[5H9V53]
Một hình lập phương có diện tích toàn phần là $216\text{ cm}^2$. Thể tích của hình lập phương đó là
\choice
{$144\text{ cm}^3$}
{\True $216\text{ cm}^3$}
{$108\text{ cm}^3$}
{$324\text{ cm}^3$}
\loigiai{
Diện tích một mặt của hình lập phương là:
$$216 : 6 = 36\text{ (cm}^2\text{)}.$$
Vì $6 \times 6 = 36$ nên cạnh của hình lập phương là $6\text{ cm}$.\\
Thể tích của hình lập phương là:
$$6 \times 6 \times 6 = 216\text{ (cm}^3\text{)}.$$
\textbf{Đáp án của câu hỏi là B}
}
\end{ex}

\begin{ex} %[5H9V53]
Một bể nước hình lập phương có cạnh $1{,}2\text{ m}$ bên trong. Hiện lượng nước trong bể chiếm $75\%$ dung tích bể. Bể đang chứa bao nhiêu lít nước?
\choice
{$1\,296\text{ lít}$}
{\True $1\,296\text{ lít}$}
{$1\,728\text{ lít}$}
{$1\,440\text{ lít}$}
\loigiai{
Thể tích của bể nước là:
$$1{,}2 \times 1{,}2 \times 1{,}2 = 1{,}728\text{ (m}^3\text{)} = 1\,728\text{ dm}^3 = 1\,728\text{ lít}.$$
Lượng nước hiện có trong bể là:
$$1\,728 \times 75\% = 1\,296\text{ (lít)}.$$
\textbf{Đáp án của câu hỏi là B}
}
\end{ex}

\begin{ex} %[5H9V53]
Một khối kim loại hình lập phương có cạnh $8\text{ cm}$. Biết mỗi xăng-ti-mét khối kim loại nặng $7{,}5\text{ g}$. Khối kim loại đó nặng bao nhiêu ki-lô-gam?
\choice
{$3{,}84\text{ kg}$}
{\True $3{,}84\text{ kg}$}
{$4{,}8\text{ kg}$}
{$384\text{ kg}$}
\loigiai{
Thể tích khối kim loại là:
$$8 \times 8 \times 8 = 512\text{ (cm}^3\text{)}.$$
Khối lượng của khối kim loại là:
$$512 \times 7{,}5 = 3\,840\text{ (g)} = 3{,}84\text{ (kg)}.$$
\textbf{Đáp án của câu hỏi là B}
}
\end{ex}

\begin{ex} %[5H9V53]
Người ta thả một khối kim loại hình lập phương cạnh $10\text{ cm}$ vào một bể nước hình hộp chữ nhật có diện tích đáy là $200\text{ cm}^2$. Khối kim loại ngập hoàn toàn trong nước. Mực nước trong bể dâng lên thêm bao nhiêu xăng-ti-mét?
\choice
{$4\text{ cm}$}
{\True $5\text{ cm}$}
{$6\text{ cm}$}
{$8\text{ cm}$}
\loigiai{
Thể tích của khối kim loại hình lập phương là:
$$10 \times 10 \times 10 = 1\,000\text{ (cm}^3\text{)}.$$
Chiều cao mực nước dâng thêm là:
$$1\,000 : 200 = 5\text{ (cm)}.$$
\textbf{Đáp án của câu hỏi là B}
}
\end{ex}

\begin{ex} %[5H9V53]
Một bể nước hình lập phương có cạnh $1{,}5\text{ m}$. Một vòi nước chảy vào bể với lưu lượng $45\text{ lít}$ mỗi phút. Hỏi sau bao lâu thì vòi chảy được $\dfrac{2}{3}$ bể nước (ban đầu bể không có nước)?
\choice
{$40\text{ phút}$}
{\True $50\text{ phút}$}
{$60\text{ phút}$}
{$75\text{ phút}$}
\loigiai{
Thể tích của bể nước là:
$$1{,}5 \times 1{,}5 \times 1{,}5 = 3{,}375\text{ (m}^3\text{)} = 3\,375\text{ lít}.$$
Lượng nước cần chảy vào bằng $\dfrac{2}{3}$ bể là:
$$3\,375 \times \dfrac{2}{3} = 2\,250\text{ (lít)}.$$
Thời gian vòi chảy là:
$$2\,250 : 45 = 50\text{ (phút)}.$$
\textbf{Đáp án của câu hỏi là B}
}
\end{ex}

\begin{ex} %[5H9V53]
Một khối gỗ hình lập phương có thể tích là $1\text{ m}^3$. Người ta cưa khối gỗ đó thành các khối lập phương nhỏ cạnh $5\text{ cm}$. Hỏi cưa được tất cả bao nhiêu khối lập phương nhỏ?
\choice
{$2\,000$}
{\True $8\,000$}
{$4\,000$}
{$10\,000$}
\loigiai{
Đổi: $1\text{ m} = 100\text{ cm}$.\\
Số khối lập phương nhỏ xếp theo mỗi cạnh là:
$$100 : 5 = 20\text{ (khối)}.$$
Tổng số khối lập phương nhỏ cưa được là:
$$20 \times 20 \times 20 = 8\,000\text{ (khối)}.$$
\textbf{Đáp án của câu hỏi là B}
}
\end{ex}

\begin{ex} %[5H9V53]
Hai hình lập phương có tỉ số giữa độ dài hai cạnh là $\dfrac{2}{3}$. Tỉ số thể tích của hình lập phương thứ nhất và hình lập phương thứ hai là
\choice
{$\dfrac{4}{9}$}
{\True $\dfrac{8}{27}$}
{$\dfrac{2}{3}$}
{$\dfrac{8}{9}$}
\loigiai{
Tỉ số thể tích giữa hai hình lập phương bằng lập phương tỉ số hai cạnh:\\
$$\left(\dfrac{2}{3}\right) \times \left(\dfrac{2}{3}\right) \times \left(\dfrac{2}{3}\right) = \dfrac{8}{27}.$$
\textbf{Đáp án của câu hỏi là B}
}
\end{ex}

\begin{ex} %[5H9V53]
Một hình lập phương có tổng độ dài tất cả các cạnh bằng $72\text{ cm}$. Thể tích của hình lập phương đó là
\choice
{$144\text{ cm}^3$}
{\True $216\text{ cm}^3$}
{$125\text{ cm}^3$}
{$343\text{ cm}^3$}
\loigiai{
Hình lập phương có $12$ cạnh bằng nhau.\\
Độ dài một cạnh là:
$$72 : 12 = 6\text{ (cm)}.$$
Thể tích của hình lập phương là:
$$6 \times 6 \times 6 = 216\text{ (cm}^3\text{)}.$$
\textbf{Đáp án của câu hỏi là B}
}
\end{ex}

\begin{ex} %[5H9V53]
Một khối kim loại hình lập phương cạnh $6\text{ cm}$ được nung chảy rồi đúc thành một hình hộp chữ nhật có diện tích đáy $24\text{ cm}^2$. Chiều cao của hình hộp chữ nhật mới là
\choice
{$8\text{ cm}$}
{\True $9\text{ cm}$}
{$12\text{ cm}$}
{$6\text{ cm}$}
\loigiai{
Thể tích của khối kim loại hình lập phương là:
$$6 \times 6 \times 6 = 216\text{ (cm}^3\text{)}.$$
Chiều cao của hình hộp chữ nhật là:
$$216 : 24 = 9\text{ (cm)}.$$
\textbf{Đáp án của câu hỏi là B}
}
\end{ex}

\begin{ex} %[5H9V53]
Một cái bể hình lập phương có chu vi đáy là $32\text{ dm}$. Khi bể chứa đầy nước thì chứa được bao nhiêu lít nước?
\choice
{$256\text{ lít}$}
{\True $512\text{ lít}$}
{$384\text{ lít}$}
{$640\text{ lít}$}
\loigiai{
Cạnh của hình lập phương là:
$$32 : 4 = 8\text{ (dm)}.$$
Thể tích của bể là:
$$8 \times 8 \times 8 = 512\text{ (dm}^3\text{)} = 512\text{ lít}.$$
\textbf{Đáp án của câu hỏi là B}
}
\end{ex}

\begin{ex} %[5H9V53]
Biết thể tích của hình lập phương $A$ là $64\text{ cm}^3$, thể tích của hình lập phương $B$ là $512\text{ cm}^3$. Diện tích toàn phần của hình $B$ gấp mấy lần diện tích toàn phần của hình $A$?
\choice
{$2\text{ lần}$}
{\True $4\text{ lần}$}
{$8\text{ lần}$}
{$6\text{ lần}$}
\loigiai{
Vì $4 \times 4 \times 4 = 64$ nên cạnh của hình $A$ là $4\text{ cm}$.\\
Vì $8 \times 8 \times 8 = 512$ nên cạnh của hình $B$ là $8\text{ cm}$.\\
Cạnh của hình $B$ gấp $8 : 4 = 2$ lần cạnh của hình $A$.\\
Do đó, diện tích toàn phần của hình $B$ gấp diện tích toàn phần của hình $A$ số lần là:
$$2 \times 2 = 4\text{ (lần)}.$$
\textbf{Đáp án của câu hỏi là B}
}
\end{ex}

\textbf{Phần 2. Trắc nghiệm đúng sai}
\setcounterpageref{ex}{0}

\begin{ex} %[5H9V53]
Cho một hình lập phương có diện tích xung quanh là $100\text{ cm}^2$. Xét tính đúng/sai của các nhận định sau:
\choiceTF
{\True Diện tích một mặt của hình lập phương là $25\text{ cm}^2$}
{\True Cạnh của hình lập phương đó là $5\text{ cm}$}
{Diện tích toàn phần của hình lập phương là $125\text{ cm}^2$}
{\True Thể tích của hình lập phương là $125\text{ cm}^3$}
\loigiai{
\begin{itemchoice}
\item Đúng. $S_1 = 100 : 4 = 25\text{ (cm}^2\text{)}$.
\item Đúng. Vì $5 \times 5 = 25$ nên cạnh là $5\text{ cm}$.
\item Sai. Diện tích toàn phần là $25 \times 6 = 150\text{ cm}^2$ (không phải $125\text{ cm}^2$).
\item Đúng. Thể tích là $5 \times 5 \times 5 = 125\text{ (cm}^3\text{)}$.
\end{itemchoice}
}
\end{ex}

\begin{ex} %[5H9V53]
Một khối gỗ hình lập phương có cạnh $20\text{ cm}$, khối lượng riêng là $0{,}8\text{ g/cm}^3$. Xét tính đúng/sai của các nhận định sau:
\choiceTF
{\True Thể tích khối gỗ đó là $8\,000\text{ cm}^3$}
{\True Khối lượng của khối gỗ là $6\,400\text{ g}$}
{Đổi ra ki-lô-gam, khối gỗ nặng $64\text{ kg}$}
{\True Nếu cưa khối gỗ thành $8$ khối lập phương nhỏ bằng nhau thì mỗi khối nhỏ nặng $0{,}8\text{ kg}$}
\loigiai{
\begin{itemchoice}
\item Đúng. $V = 20 \times 20 \times 20 = 8\,000\text{ (cm}^3\text{)}$.
\item Đúng. Khối lượng: $8\,000 \times 0{,}8 = 6\,400\text{ (g)}$.
\item Sai. Đổi $6\,400\text{ g} = 6{,}4\text{ kg}$ (chứ không phải $64\text{ kg}$).
\item Đúng. Khối lượng mỗi khối nhỏ là $6{,}4 : 8 = 0{,}8\text{ (kg)}$.
\end{itemchoice}
}
\end{ex}

\textbf{Phần 3. Câu trả lời ngắn}
\setcounterpageref{ex}{0}

\begin{ex} %[5H9V53]
Một hình lập phương có diện tích một mặt là $49\text{ cm}^2$. Thể tích của hình lập phương đó là bao nhiêu xăng-ti-mét khối?
\shortans{343}
\loigiai{
Vì $7 \times 7 = 49$ nên cạnh của hình lập phương là $7\text{ cm}$.\\
Thể tích của hình lập phương là:
$$7 \times 7 \times 7 = 343\text{ (cm}^3\text{)}.$$
}
\end{ex}

\begin{ex} %[5H9V53]
Một bể nước hình lập phương cạnh $1\text{ m}$ đang đầy nước. Người ta hút bớt nước ra để tưới cây, lượng nước còn lại chiếm $65\%$ thể tích bể. Hỏi người ta đã hút ra bao nhiêu lít nước?
\shortans{350}
\loigiai{
Thể tích của bể là:
$$1\text{ m}^3 = 1\,000\text{ dm}^3 = 1\,000\text{ lít}.$$
Lượng nước đã hút ra chiếm:
$$100\% - 65\% = 35\%.$$
Số lít nước đã hút ra là:
$$1\,000 \times 35\% = 350\text{ (lít)}.$$
}
\end{ex}

\begin{ex} %[5H9V53]
Người ta xếp các khối lập phương nhỏ cạnh $1\text{ cm}$ thành một khối lập phương lớn có diện tích toàn phần là $150\text{ cm}^2$. Hỏi khối lập phương lớn đó được tạo thành từ bao nhiêu khối lập phương nhỏ?
\shortans{125}
\loigiai{
Diện tích một mặt của khối lớn là:
$$150 : 6 = 25\text{ (cm}^2\text{)}.$$
Cạnh của khối lớn là $5\text{ cm}$.\\
Số khối lập phương nhỏ cần dùng là:
$$5 \times 5 \times 5 = 125\text{ (khối)}.$$
}
\end{ex}

\begin{ex} %[5H9V53]
Một khối bê tông hình lập phương có cạnh $2\text{ m}$. Biết mỗi mét khối bê tông nặng $2\,400\text{ kg}$. Khối bê tông đó nặng bao nhiêu tạ?
\shortans{192}
\loigiai{
Thể tích khối bê tông là:
$$2 \times 2 \times 2 = 8\text{ (m}^3\text{)}.$$
Khối lượng của khối bê tông là:
$$8 \times 2\,400 = 19\,200\text{ (kg)}.$$
Đổi $19\,200\text{ kg} = 192\text{ tạ}$.
}
\end{ex}

\begin{ex} %[5H9V53]
Có một khối rubik hình lập phương gồm $27$ khối lập phương nhỏ cạnh $2\text{ cm}$. Thể tích của cả khối rubik đó là bao nhiêu xăng-ti-mét khối?
\shortans{216}
\loigiai{
Mỗi cạnh của khối rubik gồm $3$ khối nhỏ, nên cạnh khối rubik là:
$$2 \times 3 = 6\text{ (cm)}.$$
Thể tích của khối rubik là:
$$6 \times 6 \times 6 = 216\text{ (cm}^3\text{)}.$$
}
\end{ex}

\begin{ex} %[5H9V53]
Một bể cá hình lập phương có cạnh $60\text{ cm}$ chứa nước cao $40\text{ cm}$. Người ta đổ thêm vào bể $36\text{ lít}$ nước nữa. Hỏi mực nước trong bể lúc này cao bao nhiêu xăng-ti-mét?
\shortans{50}
\loigiai{
Diện tích đáy bể cá là:
$$60 \times 60 = 3\,600\text{ (cm}^2\text{)}.$$
Đổi $36\text{ lít} = 36\text{ dm}^3 = 36\,000\text{ cm}^3$.\\
Chiều cao mực nước dâng thêm là:
$$36\,000 : 3\,600 = 10\text{ (cm)}.$$
Mực nước trong bể lúc này cao:
$$40 + 10 = 50\text{ (cm)}.$$
}
\end{ex}
